\chapter{МОДЕЛИ АДАПТИВНОГО УПРАВЛЕНИЯ (DEC-POMDP-MARL)}\label{sec:ch3}

\section{Назначение главы и место адаптивного контура в общей архитектуре}

Во второй главе была построена математическая модель распределения ресурсов, задающая структуру системы, ограничения допустимости назначения, функции полезности и стоимости, а также целевой критерий в виде социального благосостояния. Однако сама по себе эта постановка описывает прежде всего \emph{один раунд} распределения в фиксированном срезе состояния. Для реальных edge/IoT-систем этого недостаточно, поскольку решения приходится принимать многократно в условиях нестационарной нагрузки, частичной наблюдаемости, изменяющихся сетевых характеристик и возможных отказов узлов \cite{mach_2017_MEC_survey,mao_2017_MEC_perspective,ning_2024_MARL_survey,hady_2025_MARL_RAO_survey}.

Поэтому в данной главе формализуется \textbf{адаптивный контур управления}, который должен:
\begin{enumerate}
	\item учитывать последовательный, многошаговый характер распределения ресурсов;
	\item работать при частичной наблюдаемости и локальности информации;
	\item поддерживать координацию нескольких вычислительных узлов;
	\item быть совместимым с децентрализованным исполнением в онлайн-режиме;
	\item использовать целевую функцию, согласованную с критерием эффективности, введённым во второй главе.
\end{enumerate}

Для достижения этих целей задача формализуется как децентрализованный частично наблюдаемый марковский процесс принятия решений (Dec-POMDP), а в качестве вычислительного метода решения рассматривается кооперативное многоагентное обучение с подкреплением (MARL) в парадигме централизованного обучения и децентрализованного исполнения (CTDE) \cite{bernstein_2002_complexity,oliehoek_2016_DecPOMDP,amato_2024_CTDE,ning_2024_MARL_survey}. Внутри этого семейства для диссертации выбирается архитектура QMIX как компромисс между вычислительной реализуемостью, устойчивостью обучения, пригодностью к дискретным действиям и качеством координации агентов \cite{sunehag_2017_VDN,foerster_2017_COMA,lowe_2017_MADDPG,yu_2021_MAPPO,rashid_2020_QMIX,wang_2021_QPLEX}.

\section{Почему статический механизм распределения необходимо дополнить обучаемой политикой}

Даже если на каждом шаге доступен экономически корректный слой распределения, в практической системе остаются аспекты, которые трудно задать заранее в виде фиксированных параметров. К ним относятся:
\begin{itemize}
	\item нестационарные профили поступления задач;
	\item изменение пропускной способности каналов и задержек;
	\item локальные очереди и скрытые зависимости между решениями разных узлов;
	\item отказы, восстановление и временная деградация части инфраструктуры;
	\item влияние текущего решения на будущую доступность ресурсов.
\end{itemize}

Именно поэтому в современной литературе по edge orchestration и task offloading всё шире применяются методы RL и DRL: они позволяют не только реагировать на текущее состояние, но и оптимизировать \emph{долгосрочную} накопленную полезность, включая компромиссы между немедленной выгодой и будущими потерями от перегрузки \cite{qin_2024_decentralized_task,jiang_2024_task_oriented_MARL,li_2024_satellite_MARL,yao_2025_adaptive_orchestration,hady_2025_MARL_RAO_survey}. Вместе с тем одноагентные RL-подходы плохо подходят для исследуемой постановки, поскольку управление распределено между несколькими узлами, каждый из которых видит лишь локальную часть среды.

Отсюда следует, что задача должна быть представлена как кооперативная многоагентная последовательная задача принятия решений. Её естественный формализм --- Dec-POMDP.

\section{Формализация задачи адаптивного управления как Dec-POMDP}

\subsection{Базовая модель}

Задача управления распределением ресурсов формализуется в виде децентрализованного частично наблюдаемого марковского процесса принятия решений (Dec-POMDP):
\begin{equation}
	\mathcal{M}=\langle \mathcal{N},\mathcal{S},\{\mathcal{A}_j\}_{j\in\mathcal{N}},P,r,\{\mathcal{O}_j\}_{j\in\mathcal{N}},\Omega,\gamma \rangle,
	\label{eq:ch3_decpomdp}
\end{equation}
где:
\begin{itemize}
	\item $\mathcal{N}=\{1,2,\dots,k\}$ --- множество обучаемых агентов, ассоциированных с периферийными вычислительными узлами $n_j \in N$;
	\item $\mathcal{S}$ --- множество глобальных состояний среды;
	\item $\mathcal{A}_j$ --- конечное множество действий агента $j$;
	\item $P(s_{t+1}\mid s_t,\mathbf{a}_t)$ --- вероятностная функция переходов (описывает, как состояние системы меняется в результате действий всех агентов);
	\item $r(s_t,\mathbf{a}_t,s_{t+1})$ --- общее (глобальное) вознаграждение команды;
	\item $\mathcal{O}_j$ --- пространство наблюдений агента $j$;
	\item $\Omega(o_t^1,\dots,o_t^k\mid s_t,\mathbf{a}_{t-1})$ --- функция наблюдений (описывает, какое локальное наблюдение получит агент после перехода в новое состояние);
	\item $\gamma\in[0,1)$ --- коэффициент дисконтирования (параметр, определяющий вес будущих вознаграждений).
\end{itemize}

Такое представление соответствует общепринятому описанию кооперативных многоагентных задач с частичной наблюдаемостью \cite{bernstein_2002_complexity,oliehoek_2016_DecPOMDP,amato_2024_coop_intro}. В отличие от классического MDP, ни один агент не наблюдает состояние $s_t$ полностью; решение принимается по собственной истории наблюдений.

\subsection{Глобальное состояние}

Глобальное состояние $s_t \in \mathcal{S}$ описывает тот минимальный набор переменных, который необходим для адекватного моделирования дальнейшей эволюции системы. В работе предлагается использовать агрегированное состояние вида
\begin{equation}
	s_t = \Bigl(\{\mathrm{Avail}_j(t)\}_{j=1}^{k},\; \{Q_j^{\tau}(t)\}_{j=1}^{k},\; \{w_{ij}(t)\}_{i,j},\; \mathbf{z}_t\Bigr),
	\label{eq:ch3_state}
\end{equation}
где $Q_j^{\tau}(t)$ очердь задач на каждом узле $n_j$, а $\mathbf{z}_t$ --- дополнительные агрегированные признаки окружения (например, индикаторы отказов, статистики поступления задач и краткая история предыдущих назначений).

Состояние \eqref{eq:ch3_state} не предполагается непосредственно наблюдаемым агентами во время исполнения. Оно используется в фазе централизованного обучения и как теоретическая основа модели переходов.

\subsection{Локальные наблюдения и истории агентов}

Каждый агент $j$ на шаге $t$ получает локальное наблюдение
\begin{equation}
	o_t^j = \psi_j(s_t) + \varepsilon_t^j,
	\label{eq:ch3_obs}
\end{equation}
где $\psi_j(\cdot)$ --- оператор проекции глобального состояния на наблюдаемую агентом часть, а $\varepsilon_t^j$ --- шум измерения или неполнота информации. В практической постановке вектор $o_t^j$ может включать:
\begin{itemize}
	\item собственную доступную ёмкость узла и текущую загрузку;
	\item характеристики локальной очереди задач и недавно поступивших запросов;
	\item усечённые оценки задержек до соседних узлов;
	\item признаки локальных нарушений SLA, отказов и перегрузок;
	\item индикаторы предыдущего локального действия.
\end{itemize}

Поскольку одно наблюдение не восстанавливает состояние полностью, агент должен опираться на историю. Обозначим локальную историю как
\begin{equation}
	\mathcal{H}_t^j=(o_0^j,a_0^j,r_0,\dots,o_t^j).
	\label{eq:ch3_history}
\end{equation}
Именно зависимость политики от истории, а не только от текущего наблюдения, делает рекуррентные агентные сети естественным выбором для задачи с частичной наблюдаемостью \cite{oliehoek_2016_DecPOMDP,amato_2024_CTDE,rashid_2020_QMIX}.

\subsection{Действия агентов и дискретизация управления}

Для согласования с архитектурой QMIX пространство действий каждого агента делается конечным. В общем виде
\begin{equation}
	a_t^j \in \mathcal{A}_j = \{a^{(1)},a^{(2)},\dots,a^{(L)}\}.
	\label{eq:ch3_actions}
\end{equation}

Содержательно элементы $\mathcal{A}_j$ интерпретируются как локальные режимы управления: принять задачу, отложить обслуживание, отклонить, повысить приоритет критической задачи, перераспределить вычислительный бюджет между классами запросов и т.п. В экспериментальной среде такие режимы кодируются конечным числом \emph{шаблонов} локальной политики, а не непрерывными управляющими переменными. Это позволяет:
\begin{enumerate}
	\item сохранить вычислительную простоту выбора действия;
	\item обеспечить естественную совместимость с value-based MARL;
	\item отделить непрерывные параметры физической среды от дискретных управленческих решений.
\end{enumerate}

Такой приём широко используется в cooperative MARL, когда исходная задача управления имеет сложную структуру, но на уровне обучаемой политики её можно свести к выбору одного из конечного числа интерпретируемых управляющих режимов \cite{amato_2024_coop_intro,ning_2024_MARL_survey}.

\section{Целевая функция и построение вознаграждения}

\subsection{От социального благосостояния к по-шаговому вознаграждению}

Во второй главе основным критерием служила максимизация социального благосостояния $SW(t)$. Для последовательной задачи принятия решений этот критерий нужно превратить в по-шаговое вознаграждение. Базовый вариант --- использовать нормированное мгновенное социальное благосостояние:
\begin{equation}
	r_t^{\mathrm{sw}} = \frac{SW(t)}{Z_{\mathrm{sw}}},
	\label{eq:ch3_sw_reward}
\end{equation}
где $Z_{\mathrm{sw}}>0$ --- коэффициент нормировки, обеспечивающий численную устойчивость обучения.

Однако в чистом виде выражение \eqref{eq:ch3_sw_reward} может быть недостаточно чувствительным к редким, но критически важным событиям --- например, массовым нарушениям дедлайнов, перегрузке узлов или росту несбалансированности нагрузки. Поэтому в диссертации предлагается использовать составное глобальное вознаграждение
\begin{equation}
	r_t = \alpha\,\widetilde{SW}(t) - \beta\,\mathrm{Viol}(t) - \eta\,\mathrm{Drop}(t) - \zeta\,\mathrm{Imb}(t),
	\label{eq:ch3_global_reward}
\end{equation}
где:
\begin{itemize}
	\item $\widetilde{SW}(t)$ --- нормированное социальное благосостояние;
	\item $\mathrm{Viol}(t)$ --- число или доля нарушений дедлайнов/SLA;
	\item $\mathrm{Drop}(t)$ --- доля отброшенных задач;
	\item $\mathrm{Imb}(t)$ --- показатель дисбаланса загрузки между узлами;
	\item $\alpha,\beta,\eta,\zeta>0$ --- настраиваемые веса.
\end{itemize}

Такое построение согласуется с практикой reward shaping для системного управления ресурсами: глобальный критерий сохраняется, но дополняется членами, ускоряющими обучение и делающими редкие нежелательные события статистически заметными для агента \cite{yu_2021_MAPPO,yao_2025_adaptive_orchestration}. При этом важно, что все добавленные штрафы остаются подчинёнными основной логике социальной эффективности, а не заменяют её.

\subsection{Почему используется общее командное вознаграждение}

В кооперативной постановке распределения ресурсов локальные вознаграждения отдельных узлов могут породить конфликт между их текущими интересами и глобальной целью системы. Поэтому базовая модель главы~3 использует \textbf{общее командное вознаграждение}, одинаковое для всех агентов. Такой выбор:
\begin{itemize}
	\item напрямую согласует обучение с глобальной целью;
	\item упрощает интеграцию с критерием социального благосостояния;
	\item уменьшает риск того, что узел будет улучшать локальную метрику ценой деградации системы в целом.
\end{itemize}

Недостатком общего вознаграждения является проблема \emph{credit assignment}: неочевидно, какой агент внёс вклад в общий результат. Именно этот недостаток и мотивирует переход к value factorization методам, в частности к QMIX \cite{foerster_2017_COMA,sunehag_2017_VDN,rashid_2020_QMIX}.

\section{От RL к MARL: почему нужны кооперативные методы с факторизацией ценности}

\subsection{Оптимизационная задача обучения}

Пусть $\pi = (\pi_1,\dots,\pi_k)$ --- совместная политика агентов. Тогда задача кооперативного MARL состоит в максимизации ожидаемой дисконтированной суммы общих вознаграждений:
\begin{equation}
	\pi^* = \arg\max_{\pi} \; \mathbb{E}_{\pi}\left[\sum_{t=0}^{\infty} \gamma^t r_t\right].
	\label{eq:ch3_objective}
\end{equation}

Формально это прямое обобщение одноагентного RL. На практике же при наличии нескольких агентов возникают дополнительные трудности:
\begin{enumerate}
	\item \textbf{Нестационарность среды для отдельного агента:} поведение других агентов меняется по мере обучения \cite{lowe_2017_MADDPG,ning_2024_MARL_survey};
	\item \textbf{Экспоненциальный рост пространства совместных действий:} полный joint action space равен $\prod_j |\mathcal{A}_j|$;
	\item \textbf{Проблема credit assignment:} совместный результат трудно разложить на вклады агентов \cite{foerster_2017_COMA};
	\item \textbf{Частичная наблюдаемость:} политика должна опираться на локальные истории наблюдений, а не на полное состояние \cite{oliehoek_2016_DecPOMDP}.
\end{enumerate}

Следовательно, адаптивная модель должна одновременно обеспечивать координацию, масштабируемость и возможность локального исполнения без глобального контроллера.

\subsection{Парадигма CTDE}

Наиболее естественным компромиссом является парадигма centralized training with decentralized execution (CTDE). Её идея состоит в следующем \cite{amato_2024_CTDE,amato_2024_coop_intro,rashid_2020_QMIX}:
\begin{itemize}
	\item на этапе \textbf{обучения} допускается доступ к глобальному состоянию, совместным действиям и полному общему вознаграждению;
	\item на этапе \textbf{исполнения} каждый агент использует только локальную историю наблюдений и собственную локальную функцию выбора действия.
\end{itemize}

Парадигма CTDE особенно хорошо соответствует исследуемой предметной области. В симуляционной среде или офлайн-обучении можно собрать существенно более полную информацию о системе, чем доступно отдельному edge-узлу в боевом режиме. Однако в процессе реальной эксплуатации решения всё равно должны приниматься децентрализованно и быстро.

\section{Обоснование выбора QMIX как базового MARL-метода}

\subsection{Альтернативы и критерии выбора}

Для cooperative CTDE-задач литература предлагает несколько крупных семейств методов: независимое обучение, actor--critic методы с централизованным критиком, а также value factorization методы \cite{lowe_2017_MADDPG,foerster_2017_COMA,sunehag_2017_VDN,rashid_2020_QMIX,yu_2021_MAPPO,wang_2021_QPLEX}. Для исследуемой задачи ключевыми критериями являются:
\begin{enumerate}
	\item работа с дискретным пространством действий;
	\item устойчивость при частичной наблюдаемости;
	\item пригодность к децентрализованному исполнению;
	\item вычислительная масштабируемость по числу узлов;
	\item совместимость с глобальным командным вознаграждением.
\end{enumerate}

Для диссертации QMIX выбирается как \emph{базовый} адаптивный механизм, поскольку он даёт наиболее сбалансированное сочетание интерпретируемости, вычислительной простоты и силы кооперативной координации (подробнее см. Глава 1). При этом важно подчеркнуть, что выбор QMIX не означает утверждения о его универсальном превосходстве над всеми современными MARL-методами; речь идёт именно о наилучшем соответствии специфике рассматриваемой задачи.

\subsection{Почему monotonic factorization здесь уместна}

QMIX исходит из предположения, что совместную ценность действий можно представить как монотонную функцию индивидуальных полезностей агентов. Для задач распределения ресурсов это допущение во многих practically relevant режимах оказывается естественным: если при прочих равных локальное решение узла становится лучше (например, уменьшает локальную перегрузку или улучшает обслуживание важной задачи), то совместная ценность команды, как правило, не должна уменьшаться. Именно это позволяет использовать архитектуру, в которой централизованная оценка совместного действия сохраняет совместимость с локальными greedy-решениями.

Следует, однако, учитывать, что монотонность --- это \emph{ограничение выразительности}. Сценарии, где оптимальное действие одного агента зависит от контекста действий других немонотонным образом, QMIX моделирует хуже, чем более общие методы наподобие QPLEX. В диссертации это ограничение принимается осознанно ради устойчивости и воспроизводимости базовой модели; к его обсуждению работа возвращается в выводах главы и в перспективах развития метода.

\section{Архитектура модели QMIX для задачи управления ресурсами}

Архитектура QMIX-based системы управления состоит из следующих компонентов:
\begin{enumerate}
	\item Локальные Q-сети агентов (Agent Networks)
	\item Гибридная сеть смешивания (Mixing Network)
	\item Целевые сети (Target Networks)
	\item Буфер опыта (Experience Replay Buffer)
\end{enumerate}

\begin{figure}[h]
	\centering
	\includegraphics[width=\textwidth]{arch-qmix}
	\caption{Концептуальная архитектура QMIX}
	\label{img:arch-qmix}
\end{figure}

\subsection{Локальные Q-сети агентов (Agent Networks)}

Каждый агент $j$ аппроксимирует собственную локальную utility-функцию
\begin{equation}
	Q_j(\mathcal{H}_t^j,a_t^j;\theta_j),
	\label{eq:ch3_local_q}
\end{equation}
где $\mathcal{H}_t^j$ --- локальная история наблюдений, а $\theta_j$ --- параметры нейросети агента. 

Для частично наблюдаемой среды используется рекуррентная архитектура:
\begin{itemize}
	\item \textit{Входной слой}: Векторизированное локальное наблюдение $o_j$;
	\item \textit{Рекуррентный слой}: Рекуррентная нейронная сеть (Gated Recurrent Unit, GRU) аккумулирует временную информацию о скрытых состояниях среды, что особенно важно при шумных и неполных наблюдениях \cite{rashid_2020_QMIX,amato_2024_coop_intro}. Практически это означает, что агент учитывает не только текущее состояние своего узла, но и динамику очередей, сетевых колебаний и последствия предыдущих действий;
	\item \textit{Полносвязные слои}: Два полносвязных слоя для вычисления Q-значений для каждого возможного действия;
	\item \textit{Выходной слой}: Q-функция для каждой пары (наблюдение, действие).
\end{itemize}

В минимальном виде обновление скрытого состояния записывается как
\begin{equation}
	h_t^j = \mathrm{GRU}_j\bigl(h_{t-1}^j,\,[o_t^j,a_{t-1}^j]\bigr),
	\qquad
	Q_j(\mathcal{H}_t^j,\cdot)=f_j(h_t^j;\theta_j).
	\label{eq:ch3_gru_agent}
\end{equation}

\subsection{Гибридная сеть смешивания (Mixing Network)}

Центральным элементом QMIX является mixing network, которая объединяет локальные Q-значения в совместную функцию
\begin{equation}
	Q_{\mathrm{tot}}(s_t,\mathbf{a}_t)=f_{\mathrm{mix}}\bigl(s_t, Q_1,\dots,Q_k;\phi\bigr),
	\label{eq:ch3_qtot}
\end{equation}
где $\phi$ --- параметры mixing-сети. 

Важнейшее ограничение архитектуры -- \textbf{монотонность}:
\begin{equation}
	\frac{\partial Q_{\mathrm{tot}}}{\partial Q_j} \ge 0,
	\qquad \forall j \in \mathcal{N},
	\label{eq:ch3_monotonicity}
\end{equation}
то есть совместная ценность должна быть монотонной по каждому локальному Q-значению. Иными словами, увеличение любого локального Q-значения приводит к увеличению глобального Q. Именно это условие обеспечивает согласованность между централизованной оценкой совместного действия и децентрализованным выбором локальных действий \cite{rashid_2020_QMIX}.

В практической реализации веса mixing-сети генерируются гиперсетями, зависящими от глобального состояния $s_t$. Благодаря этому комбинирование локальных ценностей остаётся зависимым от текущего состояния системы, а не является фиксированной линейной суммой, как в VDN \cite{sunehag_2017_VDN,rashid_2020_QMIX}.

\subsection{Принцип локального выбора действия}

Из условия монотонности следует, что оптимальное совместное действие может быть восстановлено из локальных greedy-решений:
\begin{equation}
	a_t^{j,*}=\arg\max_{a^j\in\mathcal{A}_j} Q_j(\mathcal{H}_t^j,a^j).
	\label{eq:ch3_local_greedy}
\end{equation}

Это свойство принципиально важно для edge-среды, поскольку устраняет необходимость в онлайн-переборе полного пространства совместных действий. Каждый узел принимает решение по своей локальной сети, а координация фактически <<вшита>> в параметры, полученные на этапе централизованного обучения.

\section{Обучение QMIX-модели}

Обучение ведётся на основе \textbf{буфера опыта (Experience Replay Buffer)} $\mathcal{B}$, содержащего переходы вида
\begin{equation}
	(s_t,\mathbf{o}_t,\mathbf{a}_t,r_t,s_{t+1},\mathbf{o}_{t+1},\mathrm{done}_t).
	\label{eq:ch3_reply_buffer}
\end{equation}

Для следующего шага сначала определяются локально greedy-действия целевых сетей:
\begin{equation}
	a_{t+1}^{j,\star}=\arg\max_{a^j\in\mathcal{A}_j} Q_j(\mathcal{H}_{t+1}^j,a^j;\theta_j).
	\label{eq:ch3_target_local_action}
\end{equation}

После этого формируется \textbf{TD-цель}:
\begin{equation}
	y_t = r_t + \gamma (1-\mathrm{done}_t)\, Q_{\mathrm{tot}}^{-}\bigl(s_{t+1},\mathbf{a}_{t+1}^{\star};\theta^{-},\phi^{-}\bigr),
	\label{eq:ch3_target}
\end{equation}
где верхний индекс $^{-}$ обозначает параметры \textbf{целевых сетей (Target Networks)}.

Функция потерь имеет стандартный квадратичный вид:
\begin{equation}
	\mathcal{L}(\theta,\phi)=\mathbb{E}_{\mathcal{B}}\left[\left(y_t-Q_{\mathrm{tot}}(s_t,\mathbf{a}_t;\theta,\phi)\right)^2\right].
	\label{eq:ch3_loss}
\end{equation}

Далее Q-сети обновляются для минимизации потерь. Параметры целевых сетей обновляются реже (каждые $\tau_\text{target}$ шагов):
\begin{equation}
	\theta_{\text{target}} \leftarrow \alpha \theta + (1 - \alpha) \theta_{\text{target}}, \quad \alpha = 0.01
\end{equation}

В силу частичной наблюдаемости обучение обычно проводится не по отдельным независимым переходам, а по усечённым последовательностям (rollouts) с backpropagation through time. Это позволяет корректно обучать рекуррентные агентные сети.

\section{Парадигма CTDE в контексте edge-систем}

В нашем подходе используется специальная парадигма \textbf{CTDE}: централизованное обучение (offline) и децентрализованное исполнение (online).

\subsection{Фаза исполнения (online, децентрализованная)}

Во время исполнения каждый узел $n_j$ располагает только локальной историей $\mathcal{H}_t^j$ и собственной агентной сетью. Алгоритм его работы выглядит следующим образом:
\begin{enumerate}
	\item получить локальное наблюдение $o_t^j$;
	\item обновить скрытое состояние рекуррентной сети;
	\item вычислить $Q_j(\mathcal{H}_t^j,a)$ для всех $a\in\mathcal{A}_j$;
	\item выбрать действие по правилу \eqref{eq:ch3_local_greedy};
	\item применить действие к локальному контуру распределения ресурсов.
\end{enumerate}

Ключевое свойство этой фазы --- отсутствие потребности в глобальном состоянии и централизованном пересчёте совместной политики. Это делает модель применимой в edge-средах с ограниченной связностью и жёсткими требованиями по задержке принятия решения.

\subsection{Фаза обучения (offline, централизованная)}

На этапе обучения, напротив, допускается использование глобального состояния $s_t$, совместных действий $\mathbf{a}_t$ и общего вознаграждения $r_t$. Это может происходить:
\begin{itemize}
	\item в симуляционной среде;
	\item по журналам телеметрии и историческим трассам;
	\item в тестовом цифровом двойнике инфраструктуры.
\end{itemize}

Такая асимметрия между обучением и исполнением соответствует реальным условиям эксплуатации распределённых вычислительных систем: собрать полную информацию для постоянного онлайн-контроля трудно, но использовать её в офлайн-обучении вполне возможно \cite{amato_2024_CTDE,yao_2025_adaptive_orchestration}.

\section{Ограничения модели и границы применимости}

Несмотря на практическую привлекательность, предложенная адаптивная модель имеет ряд ограничений.

Во-первых, Dec-POMDP в общем случае остаётся вычислительно трудноразрешимой задачей; MARL выступает здесь как приближённый метод, а не как средство нахождения точного оптимума \cite{bernstein_2002_complexity,oliehoek_2016_DecPOMDP}. Следовательно, результаты обучения следует интерпретировать как приближение к эффективной политике, а не как доказательство глобальной оптимальности.

Во-вторых, QMIX опирается на ограничение монотонности \eqref{eq:ch3_monotonicity}. Если в предметной области взаимодействие агентов имеет сильно немонотонный характер, более выразительные factorization-подходы могут оказаться предпочтительнее \cite{wang_2021_QPLEX}.

В-третьих, качество обученной политики зависит от репрезентативности сценариев, использованных при обучении. Если реальные распределения нагрузки, отказов и сетевых задержек существенно расходятся с обучающей средой, необходима повторная адаптация или перенос обучения.

Тем не менее для исследуемой задачи управления ресурсами в edge-среде предложенная связка \textit{Dec-POMDP + CTDE + QMIX} представляет собой обоснованный базовый компромисс между строгостью формализации и инженерной реализуемостью.

\section{Выводы главы}

В данной главе построена модель адаптивного управления распределением ресурсов в виде кооперативного Dec-POMDP. Показано, как сущности математической постановки главы~2 --- задачи, ресурсы, сетевые характеристики, полезности, стоимости и критерий социального благосостояния --- отображаются на компоненты последовательной многоагентной модели.

Обосновано использование общей командной функции вознаграждения, согласованной с критерием социального благосостояния, и показано, что для задачи с дискретными действиями, частичной наблюдаемостью и требованием децентрализованного исполнения наиболее естественным базовым выбором является парадигма CTDE с архитектурой QMIX. Рассмотрены структура локальных агентных сетей, mixing-сети, правило обучения и особенности онлайн-исполнения.

Тем самым глава формирует адаптивный слой решения. В следующей главе этот слой будет интегрирован с аукционным механизмом распределения ресурсов, что приведёт к построению гибридного метода MA--VCG--QMIX.

